\documentclass{subfiles}
\begin{document}
    One of the fundament abilities a computer scientist must have is that of algorithms analysis,
        which can essentially be summarized as follow: let \(P\) be some algorithmically-solvable problem, 
        and let \(\A_{1}\) and \(\A_{2}\) be two algorithms, we assume to be distinct,
        that solve \(P\); how do we choose which one is better with respect to a resource?
        
    In most cases, but not exclusively, we are interested in those algorithms whose 
        time complexity is minimal. But to do so we need some tools that allows as 
        to determine the ``optimality''. One of such tools is asymptotic notation, 
        which we breafly recall in the next section.

    In some cases asymptotic notation may be missleading; in fact it can happen 
        that and algorithm \(\A_{1}\) has a worst case complexity, 
        which is the one we are interested in, that is lower the that of algorithm \(\A_{2}\),
        yet in practice \(\A_{2}\) behaves better. We discuss this kind of analysis in \Cref{Sec:1.2}.

    \subsection{Asymptotic analysis}
    \subfile{../sub/1.1 - Asymptotic analysis}

    \subsection{Amortized analysis}
    \subfile{../sub/1.2 - Amortized analysis}
\end{document}
